% !TEX encoding = UTF-8 Unicode

\documentclass[a4paper,10pt,twocolumn]{report}
\usepackage[natbibapa,nosectionbib,tocbib,numberedbib]{apacite}

\AtBeginDocument{\renewcommand{\bibname}{Literature}}

\usepackage[utf8]{inputenc}
\usepackage[T1]{fontenc}
\DeclareUnicodeCharacter{2011}{\nobreakdash-}

\usepackage{graphicx}
\usepackage{enumerate}
\usepackage{url}
\usepackage[dvipsnames]{color}
\usepackage{framed}
\definecolor{shadecolor}{named}{Melon}
\newenvironment{corona}{%
    \begin{shaded*}%
    }{%
    \end{shaded*}%
}

\usepackage[colorinlistoftodos]{todonotes}
\usepackage{pifont}
\usepackage{lmodern}
\usepackage{listings}
\lstset{
    basicstyle=\scriptsize\ttfamily,
    columns=flexible,
    breaklines=true,
    numbers=left,
    numberstyle=\tiny,
    backgroundcolor=\color[rgb]{0.85,0.90,1}
}

\let\oldquote\quote
\let\endoldquote\endquote
\renewenvironment{quote}{\footnotesize\oldquote}{\endoldquote}

\title{Big Data and Automated Content Analysis\\ (6 ECTS)\\~\\Course Manual}
\author{
    dr. Anne Kroon\\
    Graduate School of Communication\\
    University of Amsterdam\\
    a.c.kroon@uva.nl}
\date{Academic Year 2026/27\\Semester 1, block 2}

\begin{document}
\maketitle

\chapter{About this course}

This course manual contains general information, guidelines, rules, and schedules for the Research Master course Big Data \& Automated Content Analysis Part I (6 ECTS). Please read it carefully, as it contains information regarding assessments, deadlines, and grading.

\section{Course description}
``Big data'' refers to data that are more voluminous, but often also more varied, unstructured, and dynamic than data traditionally used in communication research. In Communication Science and the Social Sciences more broadly, this includes research using internet-based data sources such as social-media content, digital news archives, platform data, and public comments. This area of research is commonly described as \emph{Computational Social Science} \citep{Lazer2009} or, when focused on communication processes and content, \emph{Computational Communication Science} \citep{Shah2015}.

This course introduces the concepts, opportunities, and challenges of computational research with digital communication data. Computational approaches can make it possible to analyse large or complex collections of texts and other digital traces; they complement rather than replace established communication-science methods such as theory-driven research design, manual content analysis, and statistical inference. Throughout the course, we emphasise that automated outputs require substantive interpretation and empirical validation.

Participants learn practical strategies for acquiring, organising, preprocessing, analysing, and documenting digital data in communication contexts. The course focuses on (a) data acquisition from files and APIs, data wrangling, exploratory analysis, and reproducible workflows; and (b) automated content analysis, including text preprocessing, dictionary and Bag-of-Words approaches, unsupervised text analysis, embeddings, topic modelling, and supervised machine-learning classification. Students learn to inspect and evaluate automated outputs, compare them with appropriate benchmarks, and reflect on validity, ethics, reproducibility, and the responsible use of AI tools.

The course uses Python as its main programming environment. Students are expected to be willing to learn to write, read, adapt, and explain Python code. Previous programming experience is helpful but not required. Students without prior experience are encouraged to complete one of the many free online introductions to Python before the course begins.


\section{Goals}

Upon completion of this course, the following goals are reached:

\input{testableobjectives}

\section{Help with practical matters}

While making your first steps with programming in Python, you will probably have many questions. Nevertheless, \url{https://google.com} and \url{https://stackoverflow.com} should be your first points of contact. After all, that is how we solve many programming problems ourselves\ldots

It is also strongly encouraged to try to solve problems together with classmates. A successful technique can be \emph{pair programming}, where two colleagues work together behind one screen (for an explanation, see \url{https://en.wikipedia.org/wiki/Pair_programming}).

If you cannot find a solution independently, you can pose questions during the Thursday lab sessions. Keep the remark in Section~\ref{sec:yourcomputer} in mind.

\chapter{Rules, assessments, and grading}
\input{grading}

\section{Presence and participation}

Attendance is compulsory. Missing more than two meetings---for whatever reason---means that the course cannot be completed.

In addition to attending meetings, students are required to prepare the assigned literature and continue working on programming tasks after the lab sessions. Completing the course requires continuous self-study; attending lab sessions alone is not sufficient.

Active academic participation is also assessed through the Week 7 project-pitch, peer-feedback, and reflection component. Students are expected to prepare their pitch and provide specific, constructive feedback to peers.

\section{Staying informed}

It is your responsibility to check the means of communication used for this course---your email account and, where applicable, Canvas or other platforms used by the lecturer---regularly, which in most cases means daily.

\section{Plagiarism \& fraud}

Plagiarism is a serious academic violation. Cases in which students use material such as online sources or other sources in their written work and present this material as their own original work without citation or referencing will be reported to the Examencommissie of the Department of Communication. If the committee concludes that a student has committed plagiarism, the course cannot be completed.

General UvA regulations concerning fraud and plagiarism apply; see also \url{https://student.uva.nl/en/content/az/plagiarism-and-fraud/plagiarism-and-fraud.html}.

\subsection{Attributing code}

Plagiarism rules apply to code in the same way as they apply to text. Just as academic writing builds on the work of other scholars, code may build on solutions developed by others or found in communities such as Stack Overflow. In submitted code, it must be transparent \emph{what is your own work and what is not}.

There is no single citation style for code, but \texttt{\#comments} can be used for this purpose, for example:

\texttt{\# The following function is copied from https://stackoverflow.com/XXXXX/XXXXX}

for an almost literal copy, or:

\texttt{\# The code in this cell is inspired by https://...; I modified Y and Z}

Make sure that there is no doubt about what is yours and what is from someone else.

\subsection{Use of AI tools}

This course is designed to develop your own analytical, programming, and research skills. You remain responsible for all work submitted in your name.

\begin{itemize}
    \item The Week 4 in-class coding exercise is an individual assessment. The use of generative-AI tools is not permitted during this exercise, subject to the applicable examination arrangements.
    \item For the final project, AI tools may be used responsibly to clarify concepts, review methodologies, or help debug analyses that you have conducted yourself. They must not replace your own analytical judgement or understanding.
    \item If you make material use of AI in the final project, disclose the tool used and briefly describe how it was used. The final work must be authored, understood, checked, and defended by you.
    \item Prefer using UvA AI Chat for privacy and data protection. Be cautious with third-party tools and do not share sensitive or personal data.
    \item Normal plagiarism and fraud rules apply to AI-assisted work, code, and written text.
\end{itemize}

\subsection{Working together}

The in-class coding exercise and the final project are \emph{individual assessments}. Students may not collaborate on code, analysis, writing, or other substantive parts of these assessments unless explicitly instructed otherwise.

The Week 7 peer-feedback activity is a required and graded form of academic interaction. Students are expected to discuss project plans and provide constructive feedback to peers. Receiving or giving feedback does not permit students to share or jointly produce assessed code, analyses, or written project material.

In some cases, overlap between final projects may be unavoidable, for example when students rely on similar data sources. If in doubt, discuss this with the lecturer in advance. It is almost always possible to find a solution if the issue is discussed before submission.

For all \textbf{ungraded} assignments and exercises, collaboration is explicitly encouraged. Experience with this course has shown that students who meet regularly with classmates to discuss weekly code examples and exercises often benefit from doing so.

\section{Deadlines and handing in}

The lecturer will specify for each assessment whether it must be handed in through Canvas, FileSender, or another designated platform. The Week 4 in-class coding exercise is completed during the scheduled course meeting; students will receive practical instructions beforehand.

For the final project, submit the report as a PDF file and submit code, notebooks, data where permitted, and other required materials in the format specified by the lecturer. When multiple files are required, compress them into one \texttt{.zip} or \texttt{.tar.gz} archive. Do not email assessments directly. If FileSender is used, send files through \url{https://filesender.surf.nl/} to the lecturer's email address.

Final projects that are not handed in on time will not be graded and receive the grade 1. This rule also applies to other assessed assignments. A deadline is met only when all required files have been submitted successfully.

\begin{corona}
\section{Your computer and operating system: your responsibility}
\label{sec:yourcomputer}

Python is a language, not a program, and there are many ways of running Python code.

It is difficult to provide individual support for issues specific to your computer, operating system, or keyboard. Things may work differently on Windows, macOS, and Linux. It is your responsibility to prepare your computer before the course starts. For example, a downloaded file may be located under \texttt{/home/damian/Downloads} on Linux, \texttt{/Users/damian/Downloads} on macOS, or \texttt{C:\textbackslash\textbackslash Users\textbackslash\textbackslash Damian\textbackslash\textbackslash Downloads} on Windows.

Therefore, ensure that the following requirements are met before the course starts:

\begin{itemize}
    \item Install a recent version of Python and ensure that you can run Jupyter Notebooks. Chapter 1 of \cite{cssbook} provides instructions; you may skip the part about R. You can use Python natively or Anaconda.
    \item Ensure that you can install third-party packages. Test this by installing \texttt{shifterator} and running \texttt{import shifterator} without an error.
    \item Ensure that your keyboard produces straight quotation marks, \texttt{"}, rather than typographical quotation marks. Test this by running \texttt{print("hello world")} in a Jupyter notebook.
    \item Create a folder for this course in an easy-to-remember location, for example \texttt{/home/damian/bdaca} on Linux, \texttt{/Users/damian/bdaca} on macOS, or \texttt{C://Users//damian//bdaca} on Windows. Ensure that you can locate files in this folder from Jupyter Notebook. Preferably, do not use spaces in folder names.
\end{itemize}

If you do not succeed after trying these steps, ask classmates for help. If that does not resolve the issue, contact the lecturer as soon as possible.
\end{corona}

\chapter{Schedule and Literature}

\input{literature}
\input{schedule}

\bibliographystyle{apacite}
\bibliography{../../resources/literature}

\end{document}

